\expandafter\def\csname HistoricalNoteText337\endcsname{Sie stehen zu den vorigen in demselben Verhältniss, wie die Quadratwurzel aus der Zahl, die die Schwerkraft misst, zu der Zahl Eins.}
\par\clearpage
\SourceLeaf{332}\GaussVSourceFolio{326}{L}{ERDMAGNETISMUS UND MAGNETOMETER.}
\noindent den horizontalen Theil der erdmagnetischen Kraft ankommt, so werden wir Kürze halber jenen immer stillschweigend verstehen, wenn von der erdmagnetischen Kraft ohne den Zusatz \textit{ganze} Kraft die Rede sein wird.

Die Quadratzahl der Menge der Schwingungen einer Nadel in einer bestimmten nach Gefallen gewählten Zeit ist ein von der besondern Beschaffenheit der Nadel abhängiges Maass der Stärke des Erdmagnetismus. Das Besondere der Nadel kommt hier aber in zweierlei Rücksicht ins Spiel: einmal, insofern der Magnetismus, dessen Träger die Nadel ist, mehr oder weniger stark sein kann, zweitens, insofern die Nadel mehr oder weniger ponderable Masse, und in schwerer oder leichter zu bewegender Gestalt enthält. Die Absonderung des zweiten Theils, des Besondern der Nadel ist nun nicht schwer. Der Einfluss des Erdmagnetismus auf die in der Nadel getrennten magnetischen Flüssigkeiten bewirkt eine Drehungskraft oder ein Drehungsmoment, insofern die Nadel nicht im magnetischen Meridian ist; dies Drehungsmoment ist desto grösser, je mehr die Nadel vom magnetischen Meridian abweicht, und am grössten in der gegen den Meridian rechtwinkligen Stellung. Dies grösste Drehungsmoment wird immer stillschweigend verstanden, wenn vom Drehungsmoment schlechthin die Rede ist; es lässt sich angeben durch ein bestimmtes Gewicht, welches auf einen Hebelarm von bestimmter Länge wirkt, mithin durch eine Zahl, sobald man Gewichte und Längen, nach beliebig gewählten Einheiten, durch Zahlen ausdrückt. Nun hängt aber dieses Drehungsmoment auf eine einfache Art, welche die Dynamik lehrt, mit der Schwingungsdauer vermittelst einer durch Figur und Gewicht der Nadel bestimmten Zwischengrösse zusammen, die man ihr Trägheitsmoment nennt, und nach bekannten Regeln berechnen kann. Ist die Nadel nicht genau ein regelmässiger Körper, oder trägt sie während sie schwingt, noch sonstigen Zubehör, so bedarf es freilich zur Ausmittelung des Trägheitsmoments noch besonderer Vorkehrungen, welche hier anzugeben zu weitläufig sein würde: jedenfalls sind Mittel dazu in unsrer Gewalt. Ist nun dies Trägheitsmoment bekannt, so kann man aus der beobachteten Schwingungsdauer der Nadel auf das Drehungsmoment zurückschliessen, welches der Erdmagnetismus durch seine Einwirkung auf die magnetischen Flüssigkeiten in der Nadel hervorbringt. Übrigens ist es sehr wohl möglich, dies Drehungsmoment auch durch directe Versuche ohne beobachtete Schwingungsdauer zu bestimmen: ein eigenthümlicher dazu dienender, seit kurzem in der Göttinger Sternwarte aufgestellter Apparat zeigt sich aller nur zu
\par\clearpage
\SourceLeaf{333}\GaussVSourceFolio{327}{R}{ERDMAGNETISMUS UND MAGNETOMETER.}
\noindent wünschenden Schärfe fähig; allein für den gegenwärtigen Zweck ist es unnöthig, dabei zu verweilen.

Dieses Drehungsmoment, welches der Erdmagnetismus an einer gegebenen Nadel erzeugt, bietet uns nun eine neue Abmessungsart der Stärke der erdmagnetischen Kraft dar, oder genauer zu reden, eine neue Form der vorigen Abmessungsart, vor welcher sie den Vorzug hat, dass der eine Theil der Individualität der Nadel nunmehr abgeschieden ist. Sie bleibt von dieser Individualität nur noch insofern abhängig, als in der Nadel ein stärkerer oder schwächerer Magnetismus entwickelt sein kann, und sobald wir \textit{diesen} auf ein absolutes Maass zurückführen können, wobei das Besondere seines Trägers gar nicht mehr in Frage kommt, wird auch die Stärke des Erdmagnetismus selbst auf ein absolutes Maass zurückgeführt sein, da nur die Zahl, welche das Drehungsmoment ausdrückt, mit der Zahl, welche den Magnetismus der Nadel misst, dividirt zu werden braucht. In der That ist dann der Abmessung des Erdmagnetismus als Einheit eine solche diesem ähnlich gedachte Kraft untergelegt, deren Wirkung auf eine Einheit des Nadel-Magnetismus in einem Drehungsmoment besteht, welches durch den Druck der Gewichtseinheit auf einen Hebelarm von der Länge der Raumeinheit gemessen wird.

Man könnte versucht sein zu glauben, dass die Last, welche eine Magnetnadel zu tragen vermag, als Maassstab für die Stärke des darin entwickelten Magnetismus dienen könne. Allein eine nähere Prüfung ergibt, dass dieses Mittel für unsern Zweck ganz unbrauchbar ist. Die Bestimmung des Tragvermögens ist überhaupt keiner Schärfe fähig, indem wiederholte Versuche sehr verschiedene Resultate dafür geben können: aber, was viel wichtiger ist, dieses Tragvermögen steht mit der Grösse der Entwicklung des Magnetismus in der Nadel, in dem Sinn, wie sie hier zu verstehen ist, nemlich insofern sie das Drehungsmoment bestimmt, in gar keinem nothwendigen Zusammenhange. Bei dem Drehungsmoment kommt der Magnetismus in allen Theilen der Nadel, auf welche der Erdmagnetismus gleichmässig und in parallelen Richtungen wirkt, in Betracht: bei dem Tragvermögen hingegen hauptsächlich der ohnehin durch die Wechselwirkung des Magnetstabs und des angehängten Eisens augenblicklich modificirt werdende Magnetismus in dem der Last zunächst liegenden Ende. Zu dem hier vorliegenden Zweck sind lediglich solche Kraftwirkungen brauchbar, welche der Magnetismus aller Theile der Nadel fast gleichmässig
\par\clearpage
\SourceLeaf{334}\GaussVSourceFolio{328}{L}{ERDMAGNETISMUS UND MAGNETOMETER.}
\noindent und in fast parallelen Richtungen ausübt, also Wirkungen in beträchtlicher Entfernung.

Eine an einem bestimmten Platze befindliche Magnetnadel übt ihre magnetische Kraft in jedem Punkte des Raumes aus, in einer Stärke und Richtung, die durch die Entfernung und Lage bestimmt werden. In der Nähe ist diese Kraft stark, aber an verschiedenen Stellen sehr ungleich; in grossen Entfernungen zwar schwach, aber dann innerhalb eines mässigen Raumes an Stärke und Richtung fast gleich. Je grösser die Entfernung, desto mehr nähert sich das Gesetz der Kraft einer einfachen Regel, welche die Theorie vollständig angibt: hier dürfen wir uns auf die Betrachtung eines Falles beschränken, der für unsern Zweck hinreicht. In einer horizontalen Fläche sei N~S die festliegende Magnetnadel, deren Kraftäusserung auf eine zweite n~s an einem Faden aufgehängte hier in Frage steht: beide in solcher gegenseitigen Lage, die die Figur hinreichend erklärt.
\begin{center}
\resizebox{.50\textwidth}{!}{\begin{tikzpicture}[x=1bp,y=-1bp,line cap=round,line join=round]
\path[use as bounding box] (0,0) rectangle (1000,385);
\begin{pgfinterruptboundingbox}
\draw[line width=2.5bp] (52,65) -- (52,329);
\draw[line width=2.5bp] (672,195) -- (941,192);
\draw[line width=2.3bp] (670,202) .. controls (670,229) and (677,252) .. (690,264);
\draw[line width=2.3bp] (676,255) -- (690,264) -- (685,249);
\draw[line width=2.3bp] (937,182) .. controls (938,158) and (931,139) .. (918,128);
\draw[line width=2.3bp] (923,143) -- (918,128) -- (934,135);
\node[inner sep=0pt,outer sep=0pt] at (52,42) {\scalebox{4.2}{\ensuremath{\mathbf{N}}}};
\node[inner sep=0pt,outer sep=0pt] at (52,352) {\scalebox{4.2}{\ensuremath{\mathbf{S}}}};
\node[inner sep=0pt,outer sep=0pt] at (654,188) {\scalebox{4.2}{\ensuremath{\mathbf{n}}}};
\node[inner sep=0pt,outer sep=0pt] at (954,190) {\scalebox{4.2}{\ensuremath{\mathbf{s}}}};
\end{pgfinterruptboundingbox}
\end{tikzpicture}
}
\end{center}
Die Wirkung der erstern Nadel auf die andere wird dann in einem Bestreben, diese zu drehen, bestehen, und zwar in dem Sinn, den die Pfeile bezeichnen, wenn die Buchstaben N~n gleichnamige Pole, z.~B. die Nordpole bedeuten, mithin S~s die Südpole. Das Drehungsmoment wird ganz auf gleiche Weise durch eine Zahl verständlich gemacht, wie oben bei der Einwirkung des Erdmagnetismus auf eine frei schwebende Nadel. Die Grösse dieses Drehungsmoments hängt aber ab von der Entfernung und von der Stärke des Magnetismus in \textit{beiden} Nadeln, so dass es z.~B. bei gleicher (hinlänglich gross vorausgesetzter) Entfernung sechsmal stärker ausfällt, wenn die eine Nadel einen doppelt, die andere einen dreifach stärkern Magnetismus trüge. Mit der Entfernung hängt aber die Wirkung so zusammen, dass bei doppelter Entfernung die Wirkung nur den achten, bei dreifacher nur den siebenundzwanzigsten Theil ihres Werths bei einfacher Entfernung behält, wobei jedoch zu bemerken ist, dass dieses Gesetz nur für sehr grosse Entfernungen hinlänglich scharf, und auf kleine nicht auszudehnen ist. Da nun alle Entfernungen, nachdem für sie einmal ein Maass als Einheit gewählt ist, durch Zahlen ausgedrückt werden, so wird jenes Gesetz auch so aus-
\par\clearpage
\SourceLeaf{335}\GaussVSourceFolio{329}{R}{ERDMAGNETISMUS UND MAGNETOMETER.}
\noindent gesprochen werden können, dass das Drehungsmoment mit dem Würfel der Entfernung multiplicirt für sehr grosse Entfernungen immer gleiches Resultat gibt, welches Product man füglich das auf die Entfernungseinheit \textit{reducirte} Drehungsmoment nennen mag, ohne zu vergessen, dass nach der eben gemachten Bemerkung das in der Entfernungseinheit wirklich stattfindende Drehungsmoment, falls jene klein ist, von dem reducirten bedeutend verschieden sein kann. Dies hindert aber durchaus nicht, das reducirte Drehungsmoment zu einem Maassstabe für den Magnetismus der Nadeln zu benutzen, und \textit{den Magnetismus derjenigen Nadel als Einheit zu betrachten, welche einer andern einen eben so grossen Magnetismus tragenden in der bezeichneten Lage ein reducirtes Drehungsmoment ertheilt, welches dem Druck der Gewichtseinheit an einem Hebelarm von der Länge der Entfernungseinheit gleichkommt.} Auf diese Weise haben wir also einen völlig klaren präcisen Begriff für die Abmessung der magnetischen Kraft einer Magnetnadel gewonnen. Eine Nadel von der zweifachen Kraft wird dann einer ihr gleichmagnetisirten ein reducirtes Drehungsmoment $=4$ ertheilen u. s. w., und allgemein wird man, sobald man die Zahl für das reducirte Drehungsmoment kennt, welches eine Nadel einer ihr gleichen ertheilt, in der Quadratwurzel aus jener Zahl das absolute Maass für die Stärke des Magnetismus jeder der beiden Nadeln haben.

Es bleibt also, um die Stärke des Erdmagnetismus auf absolutes Maass zurückführen zu können, nur noch übrig, ein Verfahren anzugeben, wodurch das Drehungsmoment, welches eine Nadel einer ihr gleichen in beträchtlicher Entfernung und in der in der Figur dargestellten Lage ertheilt, mit Schärfe bestimmt werden kann. Bei einer oberflächlichen Erwägung des im Vorhergehenden absichtlich noch bei Seite gesetzten Umstandes, dass es unmöglich ist, diese so sehr schwache Wirkung der Nadel N~S auf die Nadel n~s (welche wir einstweilen genau eben so stark magnetisirt wie N~S voraussetzen wollen) für sich rein zu beobachten, da sich letztere der überall gegenwärtigen und viel stärker wirkenden erdmagnetischen Kraft nicht entziehen lässt, könnte man diese Aufgabe für sehr schwer halten: allein gerade umgekehrt wird durch diesen Umstand selbst eine leichte Lösung gegeben. Nehmen wir an, dass in unserer Figur die gerade Linie von der Mitte der Nadel N~S durch die Nadel n~s mit dem magnetischen Meridian (von Norden nach Süden zu) zusammenfalle, so wird in dieser Lage die erdmagnetische Kraft noch gar nicht auf die Nadel n~s wirken; so wie aber die
\GaussVSourceGathering{42}
\par\clearpage
\SourceLeaf{336}\GaussVSourceFolio{330}{L}{ERDMAGNETISMUS UND MAGNETOMETER.}
\noindent Drehungskraft, welche N~S auf n~s ausübt, ihr Spiel anfängt, wird n~s von ihrer ersten Lage abgelenkt werden, und in Bewegung kommen; allein je mehr sie sich in Folge dieser Bewegung von der ersten Richtung entfernt, desto stärker strebt der Erdmagnetismus, sie dahin zurückzuführen. Die Nadel macht also Schwingungen, deren Mitte aber nicht mehr die Lage im magnetischen Meridian selbst, sondern eine dagegen mehr oder weniger geneigte ist. Diese Mitte ist zugleich die Gleichgewichtslage von der Nadel n~s, welche sie annimmt, wenn die Schwingungen zur Ruhe gekommen sind. Offenbar ist ihre Richtung nichts anderes, als das Resultat der Zusammensetzung der beiden Kräfte, welche an dem Platz der Nadel n~s der Erdmagnetismus und der Magnetismus der Nadel N~S ausüben, und die unsern Voraussetzungen zufolge um einen rechten Winkel verschiedene Richtungen haben. Nach bekannten Lehren der Statik ist also das Verhältniss der Stärke dieser Kräfte, welches zugleich das Verhältniss der durch sie erzeugten Drehungsmomente ist, aus dem Ablenkungswinkel bestimmbar, d. i. aus der Ungleichheit der beiden Ruhelagen von n~s, einmal wenn beide Kräfte wirken, zweitens wenn N~S ganz entfernt ist. Hier bietet sich nun aber noch eine wichtige Bemerkung dar. Nemlich der Ablenkungswinkel der Nadel n~s ist von der Stärke ihrer Magnetisirung ganz unabhängig, da bei verstärkter Magnetisirung offenbar \textit{beide} Drehungsmomente in gleichem Verhältniss wachsen. Wir werden dadurch der sonst allerdings schwer zu erfüllenden Bedingung, dass n~s einen eben so starken Magnetismus trage, wie N~S, ganz enthoben.

Es reducirt sich also die Bestimmung der Intensität des Erdmagnetismus auf zwei Hauptgeschäfte.

I.\quad Man beobachtet die Schwingungsdauer einer Nadel N~S, und berechnet daraus das Drehungsmoment, welches der Erdmagnetismus auf diese Nadel ausübt.

II.\quad Man hängt eine zweite Nadel n~s auf, beobachtet ihre Einstellung zuerst unter dem reinen Einfluss des Erdmagnetismus, und nachher, indem N~S in beträchtlicher Entfernung, so wie es die Figur zeigt, aufgelegt ist. Aus dem Unterschied beider Stellungen oder der Ablenkung, berechnet man, welch ein Bruchtheil die Kraft der Nadel N~S von der erdmagnetischen Kraft in der gewählten Entfernung ist; ein eben so grosser Bruchtheil von dem in I. gefundenen Drehungsmoment lehrt uns das Drehungsmoment kennen, welches in jener Entfernung die Nadel N~S einer ihr gleichen ertheilen würde; dies Resultat mit dem
\par\clearpage
\SourceLeaf{337}\GaussVSourceFolio{331}{R}{ERDMAGNETISMUS UND MAGNETOMETER.}
\noindent Würfel der Entfernung multiplicirt, gibt das reducirte Drehungsmoment; die Quadratwurzel daraus die Kraft der Nadel N~S im absoluten Maasse; endlich die in I. gefundene Zahl mit dieser Quadratwurzel dividirt, gibt die Zahl für das absolute Maass des Erdmagnetismus.

Ohne mathematische Zeichen zu gebrauchen, schien diese Darstellung der Möglichkeit, die Stärke des Erdmagnetismus durch eine Zahl auszudrücken, die von der Individualität der benutzten Magnetnadeln völlig unabhängig ist, am leichtesten verständlich: bei der wirklichen Anwendung erscheint einiges in einer etwas verschiedenen Gestalt, die aber für das Wesen der Methode gleichgültig ist, auch sind dann noch manche Nebenumstände zu berücksichtigen. Nur über ein paar Umstände wollen wir hier noch einiges beifügen.

Man hat gesehen, dass die den Abmessungen untergelegten Einheiten nur in einer Entfernungseinheit und einer Gewichtseinheit bestanden. Man muss aber nicht übersehen, dass eine Gewichtsgrösse, z.~B. ein Gramm, hier nicht das Quantum ponderabler Materie bedeutete, welches diesen Namen führt, und welches überall dasselbe ist, sondern den Druck, welches dieses Quantum Materie unter dem Einfluss der Schwerkraft an dem Beobachtungsorte ausübt. Diese Schwerkraft ist aber bekanntlich an verschiedenen Orten nicht ganz gleich, und wenn wir daher den Druck eines Gramms als Gewichtseinheit wählen, so würde nach aller Strenge die Intensität des Erdmagnetismus an verschiedenen Orten nicht mit gleichem Maasse gemessen werden. Bei der grossen Schärfe, deren die Messungen gegenwärtig fähig sind, ist es billig, diesen Unterschied nicht zu vernachlässigen. Am natürlichsten ist es, ihn dadurch zu berücksichtigen, dass man die Schwerkraft selbst auf ein absolutes Maass zurückführt, indem man als ihr Maass die doppelte Fallhöhe in der gewählten Zeiteinheit, z.~B. in einer Secunde, annimmt, und den Druck durch das Produkt der Masse in die Zahl, die die Schwerkraft misst, ausdrückt. Man übersieht leicht, dass auf diese Weise andere Zahlen sowohl für die Kraft der angewandten Magnetnadel, als für die erdmagnetische Kraft hervorgehen\HistoricalNote{337}, deren Grundlagen anstatt der vorigen zwei Einheiten jetzt drei sein werden, eine Entfernungseinheit, eine Zeiteinheit und eine Masseneinheit.

\GaussVSourceGathering{42\textsuperscript{*}}
\par\clearpage
\SourceLeaf{338}\GaussVSourceFolio{332}{L}{ERDMAGNETISMUS UND MAGNETOMETER.}
Eine Hauptschwierigkeit bei Anwendung der Methode liegt noch darin, dass das obenangeführte Gesetz (die verkehrte Proportionalität der Wirkung einer Magnetnadel zu dem Würfel der Entfernung) in zulänglicher Schärfe nur für sehr grosse Entfernungen gültig ist, in welchen die Wirkungen zu klein sind, um unmittelbar mit Schärfe beobachtet werden zu können. In mässigen Entfernungen machen sich die Abweichungen von dem Gesetze schon sehr merklich: allein die Theorie lehrt, dass in diesen Abweichungen selbst wiederum Gesetzmässigkeit Statt findet, und die Mathematik gibt Mittel an die Hand, durch Combination \textit{mehrerer} in mässigen aber ungleichen Entfernungen gemachter Versuche diese Abweichungen zu erkennen, und so gut wie ganz zu eliminiren.

Immer aber dürfen, wenn diese Elimination zulässig sein soll, die Versuche nicht bei zu kleinen Entfernungen angestellt werden: die Wirkungen bleiben daher allemal vergleichungsweise nur kleine, zu deren scharfer Abmessung die früher gebräuchlichen Mittel bei weitem nicht zureichten. Gerade dieses Bedürfniss hat die Darstellung eines neuen Apparats veranlasst, der wohl am schicklichsten mit dem Namen \textit{Magnetometer} bezeichnet werden kann, da er dazu dient, alle Grössenbestimmungen sowohl in Beziehung auf die magnetische Kraft der Nadeln, als in Beziehung auf den Erdmagnetismus, wenigstens den horizontalen Theil desselben, mit einer Genauigkeit auszuführen, die der Schärfe der feinsten astronomischen Beobachtungen gleich kommt. Man bestimmt damit die Richtung des Erdmagnetismus auf eine oder ein paar Bogensecunden genau; man beobachtet Anfang und Ende einer Schwingung auf einige Hunderttheile einer Zeitsecunde sicher, also schärfer, als die Antritte der Sterne an den Fäden eines Passagen-Instruments.

Anstatt eine bereits anderwärts gegebene Beschreibung des ohnehin jetzt schon vielverbreiteten Magnetometers zu wiederholen, beschränken wir uns hier nur darauf, einige der Eigenthümlichkeiten dieses Apparats bemerklich zu machen.

Die Stellung der an einem Faden oder einem feinen Draht aufgehängten Magnetnadel und die Veränderung dieser Stellung werden nicht, wie sonst, an der Magnetnadel selbst beobachtet, sondern an dem Spiegelbilde einer in kleine Theile getheilten Scale. Der Spiegel ist an der Magnetnadel fest, also mit derselben beweglich; die Scale hingegen ist in einer beträchtlichen Entfernung davon (15 Fuss bei den Magnetometern in Göttingen) horizontal befestigt, und hinter der Scale und etwas höher befindet sich das gegen die Mitte des Spiegels ge-
\par\clearpage
\SourceLeaf{339}\GaussVSourceFolio{333}{R}{ERDMAGNETISMUS UND MAGNETOMETER.}
\noindent richtete Fernrohr, durch welches man das 30 Fuss entfernte Spiegelbild der Scale oder eines Stücks derselben sieht. Offenbar ist nun jede Veränderung der Stellung der Magnetnadel mit einer verhältnissmässigen Veränderung des Orts des Spiegelbildes verbunden, und man übersieht leicht, wie sehr die Feinheit der Beobachtung auf diese Weise gewinnt: in der That sind, wenn die Nadel einen Fuss lang ist (und grössere hat man sonst fast niemals angewandt), die Bewegungen ihrer Enden nur ein sechzigstel so gross, wie die Bewegungen des Spiegelbildes. Der Vortheil, welchen ausserdem die grosse Entfernung des Beobachters von der Magnetnadel bei der neuen Methode gewährt, ist von selbst einleuchtend, da bei der ehemaligen Art die unmittelbare Nähe des Beobachters, so wie auch der zu nächtlichen Beobachtungen nothwendigen Beleuchtungslampe mancherlei Störungen der Nadel erregen konnte.

In dem Magnetometer werden als Magnetnadeln grosse schwere Stäbe angewandt, mit so offenbarem Gewinn für die Schärfe der Beobachtungen, dass man sich jetzt nur mit Verwunderung der äusserst kleinen Nadeln erinnert, die man vordem zu den meisten magnetischen Beobachtungen zu gebrauchen pflegte. Bei der ersten Ausführung des Magnetometers wurden Stäbe von einem Pfund Gewicht angewandt; in dem magnetischen Observatorium zu Göttingen ist ein vier Pfund schwerer Stab aufgehängt, und ähnliche Stärke haben die Nadeln der meisten für andere Örter seitdem ausgeführten Apparate; das Magnetometer in der Göttinger Sternwarte hat einen fünfundzwanzig Pfund schweren Magnetstab. Je schwerer ein Magnetstab ist, desto weniger wird er von zufälligen Störungen, kleinen Luftbewegungen u.~dergl. afficirt, desto reiner stellen also seine Bewegungen den Stand der auf ihn wirkenden magnetischen Kräfte dar. Allein man darf ja nicht vergessen, dass schwere Stäbe diesen hohen Vorrang vor leichten nur dann behaupten, wenn sie auch kräftig magnetisirt sind, und dass sie ohne diese unerlässliche Bedingung nur einem Kinde in schwerer Männerrüstung gleichen würden.

Bei der Beobachtung von Schwingungszeiten bieten die schweren Stäbe noch einige besondere ungemein schätzbare Vortheile dar, namentlich dass die Dauer einer Schwingung eine beträchtliche Zeit einnimmt, und dass die Grösse des Schwingungsbogens sich sehr langsam vermindert. Die kleinen Nadeln von weniger als einem halben Loth Gewicht, deren man sich ehedem zu solchen Zwecken bediente, haben in unsern Gegenden eine Schwingungsdauer von drei bis
\par\clearpage
\SourceLeaf{340}\GaussVSourceFolio{334}{L}{ERDMAGNETISMUS UND MAGNETOMETER.}
\noindent vier Secunden; hatte man eine Beobachtung mit Schwingungsbögen von sechzig Grad angefangen, so waren diese nach einer Viertelstunde schon so klein geworden, dass man aufhören musste, und die Schwingungen selbst musste man einzeln zählen. Die vierpfündigen gut magnetisirten Stäbe machen eine Schwingung in zwanzig, der grosse fünfundzwanzigpfündige Stab eine in zweiundvierzig Secunden; fängt man auch mit Schwingungen an, die nur wenige Grade betragen, so bleiben diese doch nach vielen Stunden noch immer gross genug für die feinsten Beobachtungen. Die Beobachtung von einer oder von einigen wenigen Schwingungen gibt die Dauer immer schon mit so vieler Schärfe, dass man nachher sich entfernen kann, und nur von Zeit zu Zeit wieder einige Aufzeichnungen zu machen braucht, wobei man über die Anzahl der Schwingungen, welche die, wie eine astronomische Uhr die Zeit gleichmässig theilende, Nadel inzwischen vollendet hat, gar nicht ungewiss bleibt. Öfters hat man in der Beobachtung der Schwingungen der fünfundzwanzigpfündigen Nadel der Göttinger Sternwarte eine Unterbrechung von acht und mehrern Stunden eintreten lassen, ohne dass die Ausmittelung der Anzahl der inzwischen vollendeten Schwingungen einer Ungewissheit ausgesetzt geblieben wäre.

Eine Hauptbestimmung des Magnetometers ist die Verfolgung des Ganges der magnetischen Declination. Jedermann weiss, dass diese jetzt in ganz Europa westlich ist, und vor zweihundert Jahren östlich war. Sie ist also von Jahr zu Jahr sich anhäufenden Veränderungen unterworfen: aber sie ist auch während eines Jahres nach den Jahreszeiten ungleich; sie ist nicht einen Tag wie den andern, ja sie wechselt an einem Tage von einer Stunde zur andern. Diese sogenannten stündlichen Änderungen (die aber an feinen Apparaten schon von einer Minute zur andern, ja oft schon in kürzern Zeitfristen merklich werden) verdienen nun eine besondere Aufmerksamkeit, und eignen sich auch ganz vorzüglich zu einem eben so angenehmen als nützlichen Geschäft solcher Besitzer von Magnetometern, denen für absolute Messungen der Declination und Intensität ein angemessenes Local oder die sonstigen Zurüstungen fehlen. Bei diesen stündlichen Veränderungen der Declination hat man die regelmässigen Bewegungen von den unregelmässigen zu unterscheiden. Erstere richten sich nach der Tageszeit, und es leidet keinen Zweifel, dass der Einfluss der Sonne, wahrscheinlich insofern sie die Erde erwärmt, die Ursache davon ist. Im Allgemeinen besteht für Europa ihr Verlauf darin, dass die Nadel des Morgens, etwa um 7 oder 8 Uhr, am
\par\clearpage
\SourceLeaf{341}\GaussVSourceFolio{335}{R}{ERDMAGNETISMUS UND MAGNETOMETER.}
\noindent meisten östlich steht, oder die westliche Abweichung am kleinsten ist, dann während der Vormittagsstunden beständig zunimmt, und Nachmittags etwa um 1 oder 2 Uhr ihren grössten Stand erreicht, von welchem sie dann allmählig wieder zurückgeht, und nach Einbruch der Nacht beinahe, oder am andern Morgen ganz, ihren vorigen Stand wieder erreicht. Im hohen Sommer beträgt diese Bewegung in unsern Gegenden etwa einen Viertelsgrad oder etwas darüber; um die Zeit des kürzesten Tages kaum halb so viel. Diese regelmässigen Bewegungen folgen mithin der Stunde jedes Orts, und treten also an dem östlicher liegenden Orte wirklich früher ein, als an dem westlichern. Allein in dieselben mischen sich unregelmässige Bewegungen so sehr oft, dass vielleicht niemals jene ganz rein erscheinen: gar nicht selten sind solche unregelmässige Bewegungen, die während einer Stunde oder in noch kleinern Zeiträumen die gewöhnlichen regelmässigen des ganzen Tages weit überflügeln. Schon vor beinahe hundert Jahren hatte \textsc{Hiorter} in Upsala die Bemerkung gemacht, dass mit Nordlichtern gleichzeitig beträchtliche Bewegungen der Magnetnadel einzutreten pflegen: seitdem ist diese Erfahrung vielfach bestätigt, und man kann nicht zweifeln, dass, wenn auch nicht die Nordlichter die Ursachen der Bewegungen der Magnetnadel selbst sind, doch diejenigen (unbekannten) Ursachen, welche die Nordlichter hervorbringen, zugleich auch auf die Magnetnadel wirken, oder den Erdmagnetismus modificiren. \textsc{Arago}, ein fleissiger Beobachter der Magnetnadel, fand fast immer in Paris starke Bewegungen derselben an solchen Tagen, wo in nördlichen Gegenden Nordlichter bemerkt wurden, woraus man schliessen konnte, dass die dabei thätigen Kräfte ihre Wirkungen in grosse Entfernungen verbreiten. Helleres Licht über diese interessante Erscheinung konnte nur von verabredeten gleichzeitigen Beobachtungen an vielen von einander entfernten Orten erwartet werden, und Hr. \textsc{von Humboldt} hat seinen vielen Verdiensten um die Lehre vom Erdmagnetismus auch das beigefügt, dass er zuerst schon vor mehrern Jahren eine solche Verabredung unter den Besitzern \textsc{Gambey}scher Nadeln eingeleitet hat, wodurch jene Bemerkung schon öfters auffallend bestätigt ist.

Die Einführung der Magnetometer gab nun Gelegenheit, diese Erscheinungen mit grösster Leichtigkeit und Schärfe zu verfolgen. Schon im Laufe des Jahrs 1834 sind an vielen Orten mit ähnlichen Apparaten eine Menge gleichzeitiger Beobachtungen an verabredeten Tagen gemacht, woraus sich ergeben hat, dass nicht bloss solche grosse Bewegungen, wie die vorhin erwähnten, sondern
